\documentclass[12pt]{article}

% a template that a friend gave, it's worked well enough for me
% i have added some packages and stuff that have proved useful

\usepackage{fancyhdr}
\usepackage{tipa}
\usepackage{fontspec}
\usepackage{amsfonts}
\usepackage{enumitem}
\usepackage[margin=1in]{geometry}
\usepackage{graphicx}
\usepackage{float}
\usepackage{amsmath}
\usepackage{braket}
\usepackage{amssymb}
\usepackage{booktabs}
\usepackage{hyperref}
\usepackage{mathtools}
\usepackage{xcolor}
\usepackage{float}
\usepackage{algpseudocodex}
\usepackage{titlesec}
\usepackage{bbm}

\pagestyle{fancy}
\fancyhf{} % sets both header and footer to nothing
\lhead{Kevin Sheng}
\setmainfont{Comic Neue}
\renewcommand{\headrulewidth}{1pt}
\setlength{\headheight}{0.75in}
\setlength{\oddsidemargin}{0in}
\setlength{\evensidemargin}{0in}
\setlength{\voffset}{-.5in}
\setlength{\headsep}{10pt}
\setlength{\textwidth}{6.5in}
\setlength{\headwidth}{6.5in}
\setlength{\textheight}{8in}
\renewcommand{\headrulewidth}{0.5pt}
\renewcommand{\footrulewidth}{0.3pt}
\setlength{\textwidth}{6.5in}
\usepackage{setspace}
\usepackage{multicol}
\usepackage{float}
\setlength{\columnsep}{1cm}
\setlength\parindent{24pt}
\usepackage [english]{babel}
\usepackage [autostyle, english = american]{csquotes}
\MakeOuterQuote{"}

\setlength{\parskip}{6pt}
\setlength{\parindent}{0pt}

\titlespacing\section{0pt}{12pt plus 4pt minus 2pt}{0pt plus 2pt minus 2pt}
\titlespacing\subsection{0pt}{12pt plus 4pt minus 2pt}{0pt plus 2pt minus 2pt}
\titlespacing\subsubsection{0pt}{12pt plus 4pt minus 2pt}{0pt plus 2pt minus 2pt}

\hypersetup{colorlinks=true, urlcolor=blue}

\newcommand{\correction}[1]{\textcolor{red}{#1}}

\usepackage{fancyvrb}

\rhead{Math 168}

\DeclareMathOperator{\Fr}{Fr}
\DeclareMathOperator{\Tr}{Tr}
\newcommand{\lra}{\xLeftrightarrow}
\newcommand{\ra}{\xRightarrow}

\setlist[enumerate,1]{label=(\alph*)}

\begin{document}

\section{Newman 7.9}

I ran this code:
\begin{python}
import networkx as nx

edges = [
    (0, 1), (1, 2), (0, 2), (3, 4), (4, 5),
    (5, 6), (3, 6), (3, 5), (1, 3), (2, 6)
]
groups = [1, 1, 1, 2, 2, 2, 2]

g = nx.from_edgelist(edges)

n, m = len(g.nodes), len(g.edges)
q = 0
for i in range(n):
    for j in range(n):
        expected = g.degree[i] * g.degree[j] / (2 * m)
        q += (groups[i] == groups[j]) * (g.has_edge(i, j) - expected)
q /= 2 * m

print(f"modularity: {q}")
\end{python}
and found that the modularity was \boxed{0.28}.

\pagebreak

\section{Newman 7.16}

\subsection{Part 1}

The $e_r$ values are just one half of the fraction of edges that are within group $r$:
\begin{center}
      \begin{tabular}{|cccc|}
            \hline
            Black & Hispanic & White & Other \\
            \hline
            0.129 & 0.0785   & 0.153 & 0.008 \\
            \hline
      \end{tabular}
\end{center}
Assuming that this survey assumes no polygamy, all nodes have degree $1$ and $m=2n$.
Therefore, $a_r$ is a quarter of the fraction of nodes that are within group $r$.
Their values are as follows:
\begin{center}
      \begin{tabular}{|cccc|}
            \hline
            Black  & Hispanic & White & Other  \\
            \hline
            0.0765 & 0.0564   & 0.1   & 0.0171 \\
            \hline
      \end{tabular}
\end{center}
Combining these two values gives us a modularity of about \boxed{0.35}.

\subsection{Part 2}

We use a similar method to get these values for $e_r$:
\begin{center}
      \begin{tabular}{|ccc|}
            \hline
            Democrat & Independent & Republican \\
            \hline
            0.125    & 0.075       & 0.15       \\
            \hline
      \end{tabular}
\end{center}
and these values for $a_r$:
\begin{center}
      \begin{tabular}{|ccc|}
            \hline
            Democrat & Independent & Republican \\
            \hline
            0.08625  & 0.0625      & 0.09875    \\
            \hline
      \end{tabular}
\end{center}
which result in a modularity of \boxed{0.329}.

\subsection{Part 3}

Given that the modularity of both these approximate networks is positive,
people are certainly more likely to marry those of similar race and political
alignment.

\section{Newman 10.1++}

\begin{enumerate}
      \item The diameter of a clique is $1$.
            Every node is directly connected to every other node,
            so the longest path between any two nodes is still just one edge long.

      \item The longest path would be \boxed{2L} edges long, from the bottom left
            to the top right corner.
            If we extend this to $d$ dimensions, the longest path length becomes \boxed{dL}.
            Now, a lattice in $d$ dimensions with side length $L$ has $(L+1)^d$ nodes,
            so setting $n$ equal to that gives \boxed{dL=d\left(\sqrt[d]{n}-1\right)}.

      \item Only the first kind of network displays the small world effect,
            since the overall growth of the second one dominates logarithmic growth.
\end{enumerate}

\section{Stanford Networks}

I chose two random networks out of the first category.
\begin{enumerate}[label=\arabic*.]
      \item \textbf{LiveJournal social network}
            \begin{enumerate}
                  \item The largest WCC has $4843953$ nodes, which is around 99.9\% of all nodes in the network.
                        If that isn't a giant component, IDK what is.
                  \item The diameter is just $16$, which is actually approximately
                        equal to $\log n \approx 15.39$ in this network.
                        Yeah, I'd say it exhibits a small-world effect.
            \end{enumerate}
      \item \textbf{High-energy physics citation network}
            \begin{enumerate}
                  \item The largest WCC is $34401$ nodes, which is around 99.6\% of $n$ in this case.
                        Still giant, though.
                  \item The diameter of this network is $12$.
                        I'd say that still qualifies for the small-world effect,
                        though it may depend on one's constant factor.
            \end{enumerate}
\end{enumerate}

\end{document}
